\begin{thebibliography}{104}
\bibitem{ref1} Manindra Agrawal, Eric Allender, and Samir Datta. On $\mathrm{TC}^0$ , $\mathrm{AC}^0$ , and Arithmetic Circuits. J. Comput. Syst. Sci., 60(2):395–421, 2000. Conference version appeared in the proceedings of CCC 1997.
\bibitem{ref2} Manindra Agrawal, Chandan Saha, Ramprasad Saptharishi, and Nitin Saxena. Jacobian Hits Circuits: Hitting Sets, Lower Bounds for Depth-D Occur-k Formulas and Depth-3 Transcendence Degree-k Circuits. SIAM J. Comput., 45(4):1533–1562, 2016. Conference version appeared in the proceedings of STOC 2012.
\bibitem{ref3} Manindra Agrawal and V. Vinay. Arithmetic Circuits: A Chasm at Depth Four. In 49th Annual IEEE Symposium on Foundations of Computer Science, FOCS 2008, October 25-28, 2008, Philadelphia, PA, USA, pages 67–75. IEEE Computer Society, 2008.
\bibitem{ref4} Miklós Ajtai. $\Sigma^1_1$ -Formulae on finite structures. Annals of Pure and Applied Logic, 24(1):1–48, 1983.
\bibitem{ref5} Boris Alexeev, Michael A. Forbes, and Jacob Tsimerman. Tensor rank: Some lower and upper bounds. In Proceedings of the 26th Annual IEEE Conference on Computational Complexity, CCC 2011, San Jose, California, USA, June 8-10, 2011, pages 283–291. IEEE Computer Society, 2011.
\bibitem{ref6} Eric Allender and Michal Koucký. Amplifying lower bounds by means of self-reducibility. J. ACM, 57(3):14:1–14:36, 2010. Conference version appeared in the proceedings of CCC 2008.
\bibitem{ref7} Noga Alon and Ravi B. Boppana. The monotone circuit complexity of boolean functions. Combinatorica, 7(1):1–22, 1987.
\bibitem{ref8} Noga Alon, Mrinal Kumar, and Ben Lee Volk. Unbalancing Sets and an Almost Quadratic Lower Bound for Syntactically Multilinear Arithmetic Circuits. In Rocco A. Servedio, editor, 33rd Computational Complexity Conference, CCC 2018, June 22-24, 2018, San Diego, CA, USA, volume 102 of LIPIcs, pages 11:1–11:16. Schloss Dagstuhl - Leibniz-Zentrum fuer Informatik, 2018.
\bibitem{ref9} Matthew Anderson, Michael A. Forbes, Ramprasad Saptharishi, Amir Shpilka, and Ben Lee Volk. Identity Testing and Lower Bounds for Read-k Oblivious Algebraic Branching Programs. TOCT, 10(1):3:1–3:30, 2018. Conference version appeared in the proceedings of CCC 2016.
\bibitem{ref10} A. E. Andreev. On a method for obtaining more than quadratic effective lower bounds for thecomplexity of $\pi$-schemes. Moscow Univ. Math. Bull., 42:63–66, 1987.
\bibitem{ref11} Vikraman Arvind and Srikanth Srinivasan. On the hardness of the noncommutative determinant. Computational Complexity, 27(1):1–29, 2018. Conference version appeared in the proceedings of STOC 2010.
\bibitem{ref12} Nikhil Balaji, Nutan Limaye, and Srikanth Srinivasan. An almost cubic lower bound for $\Sigma\Pi\Sigma$ circuits computing a polynomial in VP. Electronic Colloquium on Computational Complexity (ECCC), 23:143, 2016. URL: http://eccc.hpi-web.de/report/2016/143.
\bibitem{ref13} Walter Baur and Volker Strassen. The Complexity of Partial Derivatives. Theor. Comput. Sci., 22:317–330, 1983.
\bibitem{ref14} Norbert Blum. A Boolean Function Requiring 3n Network Size. Theor. Comput. Sci., 28:337–345, 1984. doi:10.1016/0304-3975(83)90029-4.
\bibitem{ref15} Allan Borodin, Alexander A. Razborov, and Roman Smolensky. On Lower Bounds for Read-K-Times Branching Programs. Computational Complexity, 3:1–18, 1993.
\bibitem{ref16} Peter Bürgisser. Completeness and Reduction in Algebraic Complexity Theory, volume 7 of Algorithms and computation in mathematics. Springer, 2000.
\bibitem{ref17} Peter Bürgisser, Michael Clausen, and Mohammad Amin Shokrollahi. Algebraic complexity theory, volume 315 of Grundlehren der mathematischen Wissenschaften. Springer, 1997.
\bibitem{ref18} Marco L. Carmosino, Russell Impagliazzo, Shachar Lovett, and Ivan Mihajlin. Hardness Amplification for Non-Commutative Arithmetic Circuits. In Rocco A. Servedio, editor, 33rd Computational Complexity Conference, CCC 2018, June 22-24, 2018, San Diego, CA, USA, volume 102 of LIPIcs, pages 12:1–12:16. Schloss Dagstuhl - Leibniz-Zentrum fuer Informatik, 2018.
\bibitem{ref19} Prerona Chatterjee, Mrinal Kumar, Adrian She, and Ben Lee Volk. A Quadratic Lower Bound for Algebraic Branching Programs. Electronic Colloquium on Computational Complexity (ECCC), page 170, 2019. URL: https://eccc.weizmann.ac.il/report/2019/170.
\bibitem{ref20} Lijie Chen and Roei Tell. Bootstrapping results for threshold circuits "just beyond" known lower bounds. In Moses Charikar and Edith Cohen, editors, Proceedings of the 51st Annual ACM SIGACT Symposium on Theory of Computing, STOC 2019, Phoenix, AZ, USA, June 23-26, 2019, pages 34–41. ACM, 2019.
\bibitem{ref21} Xi Chen, Neeraj Kayal, and Avi Wigderson. Partial Derivatives in Arithmetic Complexity and Beyond. Foundations and Trends in Theoretical Computer Science, 6(1-2):1–138, 2011.
\bibitem{ref22} Suryajith Chillara, Christian Engels, Nutan Limaye, and Srikanth Srinivasan. A Near-Optimal Depth-Hierarchy Theorem for Small-Depth Multilinear Circuits. In Mikkel Thorup, editor, 59th IEEE Annual Symposium on Foundations of Computer Science, FOCS 2018, Paris, France, October 7-9, 2018, pages 934–945. IEEE Computer Society, 2018.
\bibitem{ref23} Suryajith Chillara, Nutan Limaye, and Srikanth Srinivasan. Small-Depth Multilinear Formula Lower Bounds for Iterated Matrix Multiplication with Applications. SIAM J. Comput., 48(1):70–92, 2019. Conference version appeared in the proceedings of STOC 2001.
\bibitem{ref24} Danny Dolev, Cynthia Dwork, Nicholas Pippenger, and Avi Wigderson. Superconcentrators, Generalizers and Generalized Connectors with Limited Depth (Preliminary Version). In David S. Johnson, Ronald Fagin, Michael L. Fredman, David Harel, Richard M. Karp, Nancy A. Lynch, Christos H. Papadimitriou, Ronald L. Rivest, Walter L. Ruzzo, and Joel I. Seiferas, editors, Proceedings of the 15th Annual ACM Symposium on Theory of Computing, 25-27 April, 1983, Boston, Massachusetts, USA, pages 42–51. ACM, 1983.
\bibitem{ref25} Zeev Dvir, Guillaume Malod, Sylvain Perifel, and Amir Yehudayoff. Separating multilinear branching programs and formulas. In Howard J. Karloff and Toniann Pitassi, editors, Proceedings of the 44th Symposium on Theory of Computing Conference, STOC 2012, New York, NY, USA, May 19 - 22, 2012, pages 615–624. ACM, 2012.
\bibitem{ref26} Paul Erdős. Beweis eines Satzes von Tschebyschef. Acta Litt. Sci. Szeged, 5:194–198, January 1932.
\bibitem{ref27} Magnus Gausdal Find, Alexander Golovnev, Edward A. Hirsch, and Alexander S. Kulikov. A Better-Than-3n Lower Bound for the Circuit Complexity of an Explicit Function. In Irit Dinur, editor, IEEE 57th Annual Symposium on Foundations of Computer Science, FOCS 2016, 9-11 October 2016, Hyatt Regency, New Brunswick, New Jersey, USA, pages 89–98. IEEE Computer Society, 2016.
\bibitem{ref28} Hervé Fournier, Nutan Limaye, Guillaume Malod, and Srikanth Srinivasan. Lower Bounds for Depth-4 Formulas Computing Iterated Matrix Multiplication. SIAM J. Comput., 44(5):1173– 1201, 2015. Conference version appeared in the proceedings of STOC 2014.
\bibitem{ref29} Merrick L. Furst, James B. Saxe, and Michael Sipser. Parity, Circuits, and the Polynomial-Time Hierarchy. Mathematical Systems Theory, 17(1):13–27, 1984. Conference version appeared in the proceedings of FOCS 1981.
\bibitem{ref30} Michelangelo Grigni and Michael Sipser. Monotone Separation of Logarithmic Space from Logarithmic Depth. J. Comput. Syst. Sci., 50(3):433–437, 1995. Conference version appeared in the proceedings of CCC 1991.
\bibitem{ref31} Dima Grigoriev and Marek Karpinski. An Exponential Lower Bound for Depth 3 Arithmetic Circuits. In Proceedings of the Thirtieth Annual ACM Symposium on the Theory of Computing, Dallas, Texas, USA, May 23-26, 1998, pages 577–582, 1998.
\bibitem{ref32} Dima Grigoriev and Alexander A. Razborov. Exponential Lower Bounds for Depth 3 Arithmetic Circuits in Algebras of Functions over Finite Fields. Appl. Algebra Eng. Commun. Comput., 10(6):465–487, 2000. Conference version appeared in the proceedings of FOCS 1998.
\bibitem{ref33} Ankit Gupta, Pritish Kamath, Neeraj Kayal, and Ramprasad Saptharishi. Approaching the Chasm at Depth Four. J. ACM, 61(6):33:1–33:16, 2014. Conference version appeared in the proceedings of CCC 2013.
\bibitem{ref34} Ankit Gupta, Pritish Kamath, Neeraj Kayal, and Ramprasad Saptharishi. Arithmetic Circuits: A Chasm at Depth 3. SIAM J. Comput., 45(3):1064–1079, 2016. Conference version appeared in the proceedings of FOCS 2013.
\bibitem{ref35} Johan Håstad. Almost Optimal Lower Bounds for Small Depth Circuits. In Proceedings of the 18th Annual ACM Symposium on Theory of Computing, May 28-30, 1986, Berkeley, California, USA, pages 6–20, 1986.
\bibitem{ref36} Johan Håstad. The Shrinkage Exponent of de Morgan Formulas is 2. SIAM J. Comput., 27(1):48–64, 1998. Conference version appeared in the proceedings of FOCS 1993.
\bibitem{ref37} William Hesse, Eric Allender, and David A. Mix Barrington. Uniform constant-depth threshold circuits for division and iterated multiplication. J. Comput. Syst. Sci., 65(4):695–716, 2002. Conference version appeared in the proceedings of CCC 2001.
\bibitem{ref38} Pavel Hrubes, Avi Wigderson, and Amir Yehudayoff. Non-commutative circuits and the sum-of-squares problem. In Leonard J. Schulman, editor, Proceedings of the 42nd ACM Symposium on Theory of Computing, STOC 2010, Cambridge, Massachusetts, USA, 5-8 June 2010, pages 667–676. ACM, 2010.
\bibitem{ref39} Pavel Hrubes and Amir Yehudayoff. Arithmetic complexity in ring extensions. Theory of Computing, 7(1):119–129, 2011.
\bibitem{ref40} Pavel Hrubes and Amir Yehudayoff. On Isoperimetric Profiles and Computational Complexity. In Ioannis Chatzigiannakis, Michael Mitzenmacher, Yuval Rabani, and Davide Sangiorgi, editors, 43rd International Colloquium on Automata, Languages, and Programming, ICALP 2016, July 11-15, 2016, Rome, Italy, volume 55 of LIPIcs, pages 89:1–89:12, 2016.
\bibitem{ref41} Russell Impagliazzo, Ramamohan Paturi, and Michael E. Saks. Size-Depth Tradeoffs for Threshold Circuits. SIAM J. Comput., 26(3):693–707, 1997. Conference version appeared in the proceedings of STOC 1993.
\bibitem{ref42} Kazuo Iwama and Hiroki Morizumi. An Explicit Lower Bound of 5n o(n) for Boolean Circuits. In Krzysztof Diks and Wojciech Rytter, editors, Mathematical Foundations of Computer Science 2002, 27th International Symposium, MFCS 2002, Warsaw, Poland, August 26-30, 2002, Proceedings, volume 2420 of Lecture Notes in Computer Science, pages 353–364. Springer, 2002.
\bibitem{ref43} Mark Jerrum and Marc Snir. Some Exact Complexity Results for Straight-Line Computations over Semirings. J. ACM, 29(3):874–897, 1982.
\bibitem{ref44} K. Kalorkoti. A Lower Bound for the Formula Size of Rational Functions. SIAM J. Comput., 14(3):678–687, 1985.
\bibitem{ref45} Daniel M. Kane and Ryan Williams. Super-linear gate and super-quadratic wire lower bounds for depth-two and depth-three threshold circuits. In Proceedings of the 48th Annual ACM SIGACT Symposium on Theory of Computing, STOC 2016, Cambridge, MA, USA, June 18-21, 2016, pages 633–643, 2016.
\bibitem{ref46} Mauricio Karchmer and Avi Wigderson. Monotone Circuits for Connectivity Require Super- Logarithmic Depth. SIAM J. Discrete Math., 3(2):255–265, 1990. Conference version appeared in the proceedings of STOC 1988.
\bibitem{ref47} Mauricio Karchmer and Avi Wigderson. On Span Programs. In Proceedings of the Eigth Annual Structure in Complexity Theory Conference, San Diego, CA, USA, May 18-21, 1993, pages 102–111. IEEE Computer Society, 1993.
\bibitem{ref48} Neeraj Kayal. An exponential lower bound for the sum of powers of bounded degree polynomials. Electronic Colloquium on Computational Complexity (ECCC), 19:81, 2012. URL: http: //eccc.hpi-web.de/report/2012/081.
\bibitem{ref49} Neeraj Kayal, Nutan Limaye, Chandan Saha, and Srikanth Srinivasan. An Exponential Lower Bound for Homogeneous Depth Four Arithmetic Formulas. SIAM J. Comput., 46(1):307–335, 2017. Conference version appeared in the proceedings of FOCS 2014.
\bibitem{ref50} Neeraj Kayal and Chandan Saha. Lower Bounds for Sums of Products of Low arity Polynomials. Electronic Colloquium on Computational Complexity (ECCC), 22:73, 2015. URL: http: //eccc.hpi-web.de/report/2015/073.
\bibitem{ref51} Neeraj Kayal and Chandan Saha. Lower Bounds for Depth-Three Arithmetic Circuits with small bottom fanin. Computational Complexity, 25(2):419–454, 2016. Conference version appeared in the proceedings of CCC 2015.
\bibitem{ref52} Neeraj Kayal, Chandan Saha, and Ramprasad Saptharishi. A super-polynomial lower bound for regular arithmetic formulas. In Symposium on Theory of Computing, STOC 2014, New York, NY, USA, May 31 - June 03, 2014, pages 146–153, 2014.
\bibitem{ref53} Neeraj Kayal, Chandan Saha, and Sébastien Tavenas. An Almost Cubic Lower Bound for Depth Three Arithmetic Circuits. In 43rd International Colloquium on Automata, Languages, and Programming, ICALP 2016, July 11-15, 2016, Rome, Italy, pages 33:1–33:15, 2016.
\bibitem{ref54} Pascal Koiran. Arithmetic circuits: The chasm at depth four gets wider. Theor. Comput. Sci., 448:56–65, 2012.
\bibitem{ref55} Alexander S. Kulikov, Olga Melanich, and Ivan Mihajlin. A 5n o(n) Lower Bound on the Circuit Size over U 2 of a Linear Boolean Function. In S. Barry Cooper, Anuj Dawar, and Benedikt Löwe, editors, How the World Computes - Turing Centenary Conference and 8th Conference on Computability in Europe, CiE 2012, Cambridge, UK, June 18-23, 2012. Proceedings, volume 7318 of Lecture Notes in Computer Science, pages 432–439. Springer, 2012.
\bibitem{ref56} Mrinal Kumar and Shubhangi Saraf. Superpolynomial lower bounds for general homogeneous depth 4 arithmetic circuits. In Automata, Languages, and Programming - 41st International Colloquium, ICALP 2014, Copenhagen, Denmark, July 8-11, 2014, Proceedings, Part I, pages 751–762, 2014.
\bibitem{ref57} Mrinal Kumar and Shubhangi Saraf. Sums of Products of Polynomials in Few Variables: Lower Bounds and Polynomial Identity Testing. In 31st Conference on Computational Complexity, CCC 2016, May 29 to June 1, 2016, Tokyo, Japan, pages 35:1–35:29, 2016.
\bibitem{ref58} Mrinal Kumar and Shubhangi Saraf. On the Power of Homogeneous Depth 4 Arithmetic Circuits. SIAM J. Comput., 46(1):336–387, 2017. Conference version appeared in the proceedings of FOCS 2014.
\bibitem{ref59} Mrinal Kumar and Ben Lee Volk. Lower Bounds for Matrix Factorization. Electronic Colloquium on Computational Complexity (ECCC), 26:47, 2019.
\bibitem{ref60} Oded Lachish and Ran Raz. Explicit lower bound of 4.5n o(n) for boolean circuits. In Jeffrey Scott Vitter, Paul G. Spirakis, and Mihalis Yannakakis, editors, Proceedings on 33rd Annual ACM Symposium on Theory of Computing, July 6-8, 2001, Heraklion, Crete, Greece, pages 399–408. ACM, 2001.
\bibitem{ref61} Shachar Lovett. Computing polynomials with few multiplications. Theory of Computing, 7(1):185–188, 2011.
\bibitem{ref62} Jacques Morgenstern. Note on a Lower Bound on the Linear Complexity of the Fast Fourier Transform. J. ACM, 20(2):305–306, 1973.
\bibitem{ref63} Cody Murray and R. Ryan Williams. Circuit lower bounds for nondeterministic quasi-polytime: an easy witness lemma for NP and NQP. In Proceedings of the 50th Annual ACM SIGACT Symposium on Theory of Computing, STOC 2018, Los Angeles, CA, USA, June 25-29, 2018, pages 890–901, 2018.
\bibitem{ref64} E. I. Nechiporuk. On a Boolean function. Doklady of the Academy of Sciences of the USSR, 164(4):765–766, 1966.
\bibitem{ref65} Noam Nisan. Lower Bounds for Non-Commutative Computation (Extended Abstract). In Cris Koutsougeras and Jeffrey Scott Vitter, editors, Proceedings of the 23rd Annual ACM Symposium on Theory of Computing, May 5-8, 1991, New Orleans, Louisiana, USA, pages 410–418. ACM, 1991.
\bibitem{ref66} Noam Nisan and Avi Wigderson. Lower Bounds on Arithmetic Circuits Via Partial Derivatives. Computational Complexity, 6(3):217–234, 1997. Conference version appeared in the proceedings of FOCS 1995.
\bibitem{ref67} Igor Carboni Oliveira and Rahul Santhanam. Hardness Magnification for Natural Problems. In Mikkel Thorup, editor, 59th IEEE Annual Symposium on Foundations of Computer Science, FOCS 2018, Paris, France, October 7-9, 2018, pages 65–76. IEEE Computer Society, 2018.
\bibitem{ref68} Pavel Pudlák. Communication in Bounded Depth Circuits. Combinatorica, 14(2):203–216, 1994.
\bibitem{ref69} Pavel Pudlák. A note on the use of determinant for proving lower bounds on the size of linear circuits. Inf. Process. Lett., 74(5-6):197–201, 2000.
\bibitem{ref70} Ran Raz. On the Complexity of Matrix Product. SIAM J. Comput., 32(5):1356–1369, 2003. Conference version appeared in the proceedings of STOC 2002.
\bibitem{ref71} Ran Raz. Separation of Multilinear Circuit and Formula Size. Theory of Computing, 2(6):121– 135, 2006. Conference version appeared in the proceedings of FOCS 2004.
\bibitem{ref72} Ran Raz. Multi-linear formulas for permanent and determinant are of super-polynomial size. J. ACM, 56(2):8:1–8:17, 2009. Conference version appeared in the proceedings of STOC 2004.
\bibitem{ref73} Ran Raz. Elusive Functions and Lower Bounds for Arithmetic Circuits. Theory of Computing, 6(1):135–177, 2010. Conference version appeared in the proceedings of STOC 2008.
\bibitem{ref74} Ran Raz. Tensor-Rank and Lower Bounds for Arithmetic Formulas. J. ACM, 60(6):40:1–40:15, 2013. Conference version appeared in the proceedings of STOC 2010.
\bibitem{ref75} Ran Raz and Pierre McKenzie. Separation of the Monotone NC Hierarchy. Combinatorica, 19(3):403–435, 1999. Conference version appeared in the proceedings of FOCS 1997.
\bibitem{ref76} Ran Raz and Amir Shpilka. Lower Bounds for Matrix Product in Bounded Depth Circuits with Arbitrary Gates. SIAM J. Comput., 32(2):488–513, 2003. Conference version appeared in the proceedings of STOC 2001.
\bibitem{ref77} Ran Raz, Amir Shpilka, and Amir Yehudayoff. A Lower Bound for the Size of Syntactically Multilinear Arithmetic Circuits. SIAM J. Comput., 38(4):1624–1647, 2008. Conference version appeared in the proceedings of FOCS 2007.
\bibitem{ref78} Ran Raz and Amir Yehudayoff. Balancing Syntactically Multilinear Arithmetic Circuits. Computational Complexity, 17(4):515–535, 2008.
\bibitem{ref79} Ran Raz and Amir Yehudayoff. Lower Bounds and Separations for Constant Depth Multilinear Circuits. Computational Complexity, 18(2):171–207, 2009. Conference version appeared in the proceedings of CCC 2008.
\bibitem{ref80} A. A. Razborov. Lower bounds on monotone complexity of the logical permanent. Mathematical notes of the Academy of Sciences of the USSR, 37(6):485–493, June 1985.
\bibitem{ref81} Alexander A. Razborov. Lower bounds on the monotone complexity of some Boolean functions. Soviet Mathematics Doklady, 31:354–357, 1985.
\bibitem{ref82} John H. Reif and Stephen R. Tate. On Threshold Circuits and Polynomial Computation. SIAM J. Comput., 21(5):896–908, 1992.
\bibitem{ref83} Robert Robere, Toniann Pitassi, Benjamin Rossman, and Stephen A. Cook. Exponential Lower Bounds for Monotone Span Programs. In Irit Dinur, editor, IEEE 57th Annual Symposium on Foundations of Computer Science, FOCS 2016, 9-11 October 2016, Hyatt Regency, New Brunswick, New Jersey, USA, pages 406–415. IEEE Computer Society, 2016.
\bibitem{ref84} Eli Shamir and Marc Snir. Lower bounds on the number of multiplications and the number of additions in monotone computations. Technical report, IBM RC 6757, 1977.
\bibitem{ref85} Abhijat Sharma. An Improved Lower Bound for Depth Four Arithmetic Circuits. Master’s thesis, Indian Institute of Science, Bangalore, India, 2017.
\bibitem{ref86} Victor Shoup and Roman Smolensky. Lower Bounds for Polynomial Evaluation and Interpolation Problems. Computational Complexity, 6(4):301–311, 1997. Conference version appeared in the proceedings of FOCS 1991.
\bibitem{ref87} Amir Shpilka and Avi Wigderson. Depth-3 arithmetic circuits over fields of characteristic zero. Computational Complexity, 10(1):1–27, 2001. Conference version appeared in the proceedings of CCC 1999.
\bibitem{ref88} Amir Shpilka and Amir Yehudayoff. Arithmetic Circuits: A survey of recent results and open questions. Foundations and Trends in Theoretical Computer Science, 5(3-4):207–388, 2010.
\bibitem{ref89} Srikanth Srinivasan. Strongly Exponential Separation Between Monotone VP and Monotone VNP. Electronic Colloquium on Computational Complexity (ECCC), 26:32, 2019.
\bibitem{ref90} Volker Strassen. Die berechnungskomplexiät von elementarysymmetrischen funktionen und von iterpolationskoeffizienten. Numerische Mathematik, 20:238–251, 1973.
\bibitem{ref91} Volker Strassen. Vermeidung von divisionen. The Journal für die Reine und Angewandte Mathematik, 264:182–202, 1973.
\bibitem{ref92} Éva Tardos. The gap between monotone and non-monotone circuit complexity is exponential. Combinatorica, 8(1):141–142, 1988.
\bibitem{ref93} Sébastien Tavenas. Improved bounds for reduction to depth 4 and depth 3. Inf. Comput., 240:2–11, 2015. Conference version appeared in the proceedings of MFCS 2013.
\bibitem{ref94} Leslie G. Valiant. On Non-linear Lower Bounds in Computational Complexity. In William C. Rounds, Nancy Martin, Jack W. Carlyle, and Michael A. Harrison, editors, Proceedings of the 7th Annual ACM Symposium on Theory of Computing, May 5-7, 1975, Albuquerque, New Mexico, USA, pages 45–53. ACM, 1975.
\bibitem{ref95} Leslie G. Valiant. Graph-Theoretic Arguments in Low-Level Complexity. In Mathematical Foundations of Computer Science 1977, 6th Symposium, Tatranska Lomnica, Czechoslovakia, September 5-9, 1977, Proceedings, pages 162–176, 1977.
\bibitem{ref96} Leslie G. Valiant. Completeness Classes in Algebra. In Proceedings of the 11h Annual ACM Symposium on Theory of Computing, April 30 - May 2, 1979, Atlanta, Georgia, USA, pages 249–261, 1979.
\bibitem{ref97} Leslie G. Valiant. Negation can be exponentially powerful. Theor. Comput. Sci., 12:303–314, 1980. Conference version appeared in the proceedings of STOC 1979.
\bibitem{ref98} Leslie G. Valiant, Sven Skyum, S. Berkowitz, and Charles Rackoff. Fast Parallel Computation of Polynomials Using Few Processors. SIAM J. Comput., 12(4):641–644, 1983.
\bibitem{ref99} Vijay V. Vazirani. Approximation algorithms. Springer, 2001. URL: http://www.springer. com/computer/theoretical+computer+science/book/978-3-540-65367-7.
\bibitem{ref100} Joachim von zur Gathen and Gadiel Seroussi. Boolean Circuits Versus Arithmetic Circuits. Inf. Comput., 91(1):142–154, 1991.
\bibitem{ref101} Ryan Williams. Nonuniform ACC Circuit Lower Bounds. J. ACM, 61(1):2:1–2:32, 2014. Conference version appeared in the proceedings of CCC 2011.
\bibitem{ref102} David P. Williamson and David B. Shmoys. The Design of Approximation Algorithms. Cambridge University Press, 2011. URL: http://www.cambridge.org/de/knowledge/isbn/ item5759340/?site\_locale=de\_DE.
\bibitem{ref103} Morris Yau. Almost cubic bound for depth three circuits in VP. Electronic Colloquium on Computational Complexity (ECCC), 23:187, 2016. URL: http://eccc.hpi-web.de/report/ 2016/187.
\bibitem{ref104} Amir Yehudayoff. Separating monotone VP and VNP. In Moses Charikar and Edith Cohen, editors, Proceedings of the 51st Annual ACM SIGACT Symposium on Theory of Computing, STOC 2019, Phoenix, AZ, USA, June 23-26, 2019, pages 425–429. ACM, 2019.
\end{thebibliography}
